\section[Desenvolvimento do Projeto]{Desenvolvimento do Projeto}

\subsection{Repositório do Código Fonte do Projeto}

O projeto possui repositórios separados para o backend, frontend e wiki:

\begin{itemize}
  \item \textbf{Backend:} \url{https://tools.ages.pucrs.br/sistema-amores-fati/sistema-amores-fati-backend.git}
  \item \textbf{Frontend:} \url{https://tools.ages.pucrs.br/sistema-amores-fati/sistema-amores-fati.git}
  \item \textbf{Wiki:} \url{https://tools.ages.pucrs.br/sistema-amores-fati/wiki.git}
\end{itemize}

\subsection{Banco de Dados Utilizado}

O modelo de dados foi construído utilizando a ferramenta \textbf{dbdiagram.io}, que permite
a definição do esquema em \ac{sql} e a visualização do diagrama \ac{uml}. O diagrama pode ser
acessado em:

\url{https://dbdiagram.io/d/Amores-Fati-69bc7eeefb2db18e3bc293ee}

\subsection{Arquitetura Utilizada}

O sistema adota uma arquitetura web com separação clara entre frontend e backend,
hospedagem na \textbf{Vercel} e banco de dados relacional. A escolha priorizou
tecnologias amplamente reconhecidas pelo mercado, com foco em desempenho,
segurança, escalabilidade e facilidade de manutenção.

\subsection{Protótipos das Telas Desenvolvidas}

Os protótipos foram desenvolvidos no Figma, com padrão visual definido pelo designer
do time (cores, tipografia, botões e ícones). As telas foram validadas com os clientes
ao final da Sprint 0. O protótipo pode ser acessado em:

\url{https://www.figma.com/design/pcQ9rEVhEWim4YtQj328CC/Amores-Fati?node-id=9-251}

\subsection{Tecnologias Utilizadas}

As tecnologias utilizadas no projeto serão definidas ao longo do desenvolvimento,
priorizando soluções amplamente reconhecidas pelo mercado de trabalho em tecnologia,
com foco em desempenho, segurança e escalabilidade.
